\documentclass[12pt]{article}
\usepackage[spanish]{babel}
\usepackage{amsmath,amssymb}
\usepackage{geometry}
\geometry{margin=2.5cm}

\title{Volúmenes de Sólidos de Revolución}
\author{}
\date{}

\begin{document}
\maketitle

\section*{Ejercicio 31}
Calcular el volumen de un cono de radio $r$ y altura $h$.

Un cono es un sólido geométrico que puede obtenerse como sólido de revolución. Su volumen es una fórmula conocida en geometría espacial y puede deducirse mediante integración o comparación con un cilindro.

El volumen de un cono está dado por:
\[
V = \frac{1}{3}\pi r^2 h
\]

Este resultado indica que el volumen del cono corresponde a la tercera parte del volumen de un cilindro de igual base y altura.

\section*{Ejercicio 32}
Calcular el volumen generado por una elipse horizontal de semieje mayor $a$ y semieje menor $b$ al rotar alrededor del eje $x$.

La ecuación de la elipse centrada en el origen es:
\[
\frac{x^2}{a^2} + \frac{y^2}{b^2} = 1
\]

Despejando $y$ en función de $x$, se obtiene:
\[
y = b\sqrt{1 - \frac{x^2}{a^2}}
\]

Al girar la elipse alrededor del eje $x$, se utiliza el método de discos, ya que las secciones transversales perpendiculares al eje de rotación son circulares.

El volumen se calcula mediante:
\[
V = \pi \int_{-a}^{a} y^2 \, dx
\]

Sustituyendo:
\[
V = \pi \int_{-a}^{a} b^2 \left(1 - \frac{x^2}{a^2}\right) dx
\]

Resolviendo la integral se obtiene:
\[
V = \frac{4}{3}\pi a b^2
\]

\section*{Ejercicio 33}
Hallar el volumen que se genera al rotar el área comprendida entre la parábola
\[
y = 4x - x^2
\]
y el eje $x$, con respecto a la recta $y = 6$.

Primero se determinan los puntos de intersección con el eje $x$:
\[
4x - x^2 = 0 \Rightarrow x = 0, \, 4
\]

Como el eje de rotación es la recta horizontal $y = 6$, se emplea el método de arandelas. El radio exterior es la distancia desde $y=6$ hasta el eje $x$, y el radio interior es la distancia desde $y=6$ hasta la curva.

\[
R = 6, \qquad r = 6 - (4x - x^2)
\]

El volumen viene dado por:
\[
V = \pi \int_{0}^{4} \left(R^2 - r^2\right) dx
\]

\[
V = \pi \int_{0}^{4} \left[36 - (6 - 4x + x^2)^2\right] dx
\]

Resolviendo la integral:
\[
V = \frac{1408}{15}\pi
\]

\section*{Ejercicio 34}
Encontrar el volumen del sólido que se genera al girar la región comprendida entre las curvas
\[
y = e^x, \qquad y = 1,
\]
y la recta $x = 1$ alrededor del eje $x$.

Las curvas $y=e^x$ y $y=1$ se intersectan en $x=0$, por lo que el intervalo de integración es $[0,1]$.

Al girar alrededor del eje $x$, se utiliza el método de arandelas, donde:
\[
R = e^x, \qquad r = 1
\]

El volumen se calcula como:
\[
V = \pi \int_{0}^{1} \left(e^{2x} - 1\right) dx
\]

Integrando:
\[
V = \pi \left[ \frac{e^{2x}}{2} - x \right]_{0}^{1}
\]

Finalmente:
\[
V = \frac{\pi}{2}(e^2 - 3)
\]

\section*{Ejercicio 35}
Hallar el volumen del toroide que se genera al girar la curva
\[
(x - 3)^2 + y^2 = 4
\]
alrededor del eje $x$.

La ecuación representa un círculo de radio $r = 2$ cuyo centro se encuentra a una distancia $R = 3$ del eje de rotación. Al girar este círculo alrededor del eje $x$ se genera un toroide.

El volumen de un toroide está dado por la fórmula:
\[
V = 2\pi^2 R r^2
\]

Sustituyendo los valores:
\[
V = 2\pi^2 (3)(2^2) = 24\pi^2
\]

\section*{Ejercicio 36}
Hallar el volumen del sólido generado al girar la región comprendida entre las curvas
\[
y = 8x^2 \quad \text{y} \quad y = \sqrt{8x}
\]
alrededor del eje $x$.

Para hallar los límites de integración se igualan ambas funciones:
\[
8x^2 = \sqrt{8x}
\]

De esta ecuación se obtienen los puntos de intersección:
\[
x = 0, \quad x = \frac{1}{2}
\]

Al girar alrededor del eje $x$, se aplica el método de arandelas, donde:
\[
R = \sqrt{8x}, \qquad r = 8x^2
\]

El volumen se calcula como:
\[
V = \pi \int_{0}^{1/2} \left(8x - 64x^4\right) dx
\]

Resolviendo la integral:
\[
V = \frac{3\pi}{5}
\]

\section*{Ejercicio 37}
Usar el método más apropiado para calcular el volumen que se genera al girar el área limitada por las curvas
\[
y = 2 - x^4 \quad \text{y} \quad y = 1
\]
alrededor del eje $x$.

Los puntos de intersección se obtienen resolviendo:
\[
2 - x^4 = 1 \Rightarrow x = \pm 1
\]

Como el eje de giro es el eje $x$, se emplea el método de arandelas, donde:
\[
R = 2 - x^4, \qquad r = 1
\]

El volumen es:
\[
V = \pi \int_{-1}^{1} \left[(2 - x^4)^2 - 1\right] dx
\]

\[
V = \pi \int_{-1}^{1} (3 - 4x^4 + x^8) dx
\]

Resolviendo:
\[
V = \frac{128\pi}{45}
\]

\end{document}
